% Time-stamp: <2025-05-15 01:26:31 okumura>
\documentclass{jarticle}

\usepackage{tam}
\usepackage{fleqn}
\usepackage{url}
\usepackage{bm}
\usepackage{amsmath} 
\usepackage[dvipdfmx]{graphicx}
\usepackage[dvipdfmx]{color}
%\usepackage{color}
\usepackage{colordef}
\usepackage{amssymb}

%%%%%%%%%% Command alias %%%%%%%%%%
\newcommand{\IB}{$\bullet$}
\newcommand{\IC}{$\circ$}
\newcommand{\IS}{$*$}
%%%%%%%%%%%%%%%%%%%%%%%%%%%%%%%%%%%

\title{産業用の買電価格について}
\author{神戸大学大学院 奥村侑真}
\date{}

\def\baselinestretch{1.02}
\renewcommand{\thefootnote}{\roman{footnote}}

\begin{document}
\maketitle

\section{はじめに}
産業用の買電価格の例として，四国電力の買電価格を調べた．四国電力の料金メニューは低圧(100Vや200V)で電気を使う場合と，高圧(6000V以上)で電気を使う場合で買電価格が異なる．さらに，高圧で電気を使う場合は，(1)事務所ビル・商業施設などで電気を使用する場合，(2)工場などで電気を使用する場合の2つで価格が異なる．本研究の対象はテクシード石井工場であるため，四国電力の料金メニューのうち，高圧かつ工場使用を前提とした料金メニューに基づいて買電価格をまとめた．なお，特に記述がない限り，価格は2024年4月時点のものである．本稿の内容は，四国電力株式会社の公式ウェブサイト（\url{https://www.yonden.co.jp/business/price/plan/10.html}，2025年5月22日アクセス）を参照して作成している．

\section{買電価格の算出法}
買電価格は以下のようにして算出される．\\
電気料金=基本料金+電力量料金+再エネ賦課金\\
以下，基本料金・電力量料金，再エネ賦課金の3点について具体的に説明する．

\subsection{基本料金}
基本料金は毎月発生する料金で，契約電力に依存する．基本料金は以下のようにして算出される．\\
基本料金=基本料金単価×契約電力±力率割引(割増)\\
\subsubsection{基本料金単価}
基本料金単価は契約電力が500kW未満の場合は1412.61円/kWであり，500kW以上である場合，1935.37円/kWである．ただし，契約電力が500kW以上である場合は，500kWを上回った電力分だけに1935.37円/kWが適応されるわけではない．つまり，基本料金単価は契約電力の増加に伴って連続的に増加するわけではなく，500kWの部分が非連続となる．
なお，500kW未満の契約電力で契約しているにも関わらず，契約電力が500kW以上となってしまった場合は，四国電力との協議によってすみやかに需給契約を改めることになる。
\footnote{四国電力「業務用電力 主契約料金条件 [契約電力500キロワット未満]」，\url{https://www.yonden.co.jp/business/assets/pdf/price/plan/10/20240401_1.pdf}（2025年5月22日アクセス）}
\subsubsection{契約電力}
原則として，その月の最大需要電力と前11ヶ月の最大需要電力のうち，大きい値の方を契約電力とする．なお，需要電力とは，30分ごとの平均の消費電力のことを指す\footnote{四国電力「契約電力の決定方法について」，\url{https://www.yonden.co.jp/business/price/determine.html}（2025年5月14日アクセス）}．
\subsubsection{力率割引(割増)}
力率は，その1月のうち毎日午前8時から午後10時までの時間における平均力率とする．力率が，85\%を上回る場合は，その上回る1\%につき，基本料金を1\%割引し，85\%を下回る場合は，その下回る1\%につき，基本料金を1\%割増しとする．なお，まったく電気を使用しないその1月の力率は，85\%とみなす\footnote{四国電力「高圧電力A 主契約料金条件」，\url{https://www.yonden.co.jp/business/assets/pdf/price/plan/10/20240401_4.pdf}（2025年5月14日アクセス）}.

\subsection{電力量料金}
電力量料金は使用電力量1kWhあたりの料金であり，多くの場合，季節によって単価が異なる．電力量料金は以下のようにして算出される．\\
電力量料金=(電力量料金単価+燃料費調整単価)×使用電力量
\subsubsection{電力量料金単価} 
電力量料金単価は夏季(7〜9月)とその他季で異なる．\\
(1)夏季\\
契約電力が500kW未満の場合は29.76円/kWhであり，500kW以上である場合，26.82円/kWhである．\\
(2)その他季\\
契約電力が500kW未満である場合は28.47円/kWhであり，500kW以上である場合，25.77円/kWhである．
\subsubsection{燃料費調整単価}
発電に使用する原油，液化天然ガス，石炭といった燃料の価格変動を，電気料金に反映させた単価である．2025年2月時点では，燃料費調整単価は-7.52円である．

\subsection{再エネ賦課金}
再生可能エネルギーで発電された電気を，電力会社が一定期間・固定価格で買い取る制度（FIT制度）の費用を賄うために，電気を使用するすべての顧客が負担するものである．
再エネ賦課金は以下のようにして算出される．\\
再エネ賦課金=再エネ賦課金単価×使用電力量\\
2024年5月分から2025年度4月分の再エネ賦課金単価は3.49円/kWhである\footnote{四国電力「再生可能エネルギー発電促進賦課金・単価表」，\url{https://www.yonden.co.jp/extra/renewable.html}（2025年6月17日アクセス）}.

\section{買電料金の計算例}
2024年12月16日〜20日のテクシードの計測データをもとに，12月16日〜20日における30分ごとの平均購入電力量のうち最大の値である78kWを契約電力，12月16日〜20日で購入した電力量を4倍した4216kWhを2024年12月の使用電力量とし，力率割引を無視してテクシードの2024年12月における買電価格を計算すると，\\
基本料金=(1412.61×78)円\\
電力量料金= ((28.47-7.52)×4216)円\\
再エネ賦課金=(3.49×4216)円\\
であることより，1412.61×78+(28.47-67.52+3.49)×4216=1412.61×78+24.44×4216≒213223円となる．

\section{買電料金を目的関数に組み込む}
購買料金は，契約電力と使用電力に注目すると，契約電力が500kW未満かつその他季の場合，\\
購買料金=1412.61×契約電力 + 24.44×使用電力\\
と表すことができるため，これを目的関数に組み込むことを来週以降考える．


\end{document}
